\documentclass[10pt,a4paper,twocolumn]{article}
\usepackage[utf8]{inputenc}
\usepackage[T1]{fontenc}
\usepackage[british]{babel}
\usepackage[margin=1.5cm]{geometry}
\setlength{\columnsep}{0.8cm}
\setcounter{topnumber}{3}
\setcounter{totalnumber}{4}
\renewcommand{\topfraction}{0.9}
\renewcommand{\textfraction}{0.08}
\renewcommand{\floatpagefraction}{0.8}
\usepackage{booktabs}
\usepackage{microtype}
\usepackage{titlesec}
\usepackage[hidelinks]{hyperref}
\titlespacing*{\section}{0pt}{0.75em}{0.3em}
\titlespacing*{\subsection}{0pt}{0.6em}{0.2em}
\setlength{\parskip}{0.28em}
\setlength{\parindent}{0pt}
\title{\vspace{-1.0em}\textbf{What a domain prior can and cannot replace}\vspace{-0.4em}}
\author{\textbf{Cajetan Songwae} \\
  \small{\href{mailto:songwae2000@gmail.com}{songwae2000@gmail.com}} \\
  \small{TGDC case study, research track}}
\date{}
\begin{document}
\twocolumn[
  \begin{@twocolumnfalse}
  \maketitle
  \vspace{-2.0em}
  \end{@twocolumnfalse}
]

Companies that sell generated data build it from a description of a business. An
org chart, a list of the work it does, and rules for how that work behaves. If
the rules can supply the behaviour, you can build a client's data without ever
seeing their records.

I wanted to know whether that holds. So I picked one thing a generator has to get
right, and I tested it. How long a job takes.

\section*{Hypothesis}

I wrote my prediction down before I ran anything.

Rules reasoned from the nature of the work should capture a good share of what
real records tell you about duration. In numbers, a prior written without seeing
the target data should reach an AUC of 0.60 to 0.72. Below 0.55 would mean the
reasoning adds nothing.

If I am right, a list of job types is enough to build a generator. If I am wrong,
the behaviour has to come from the client's own records.

\section*{Method}

I used public 311 service request feeds from four cities. These are real
operational records. The people doing the work wrote them. Each feed gives the
work type, the owning department, how the request came in, and when it opened and
closed. I took two full calendar weeks from each city, the weeks of 4 May and 8
June 2026.

\begin{table*}[t]
\centering
\footnotesize
\begin{tabular}{lrrrr}
\toprule
city & train & test & median resolution & categories \\
\midrule
New York & 72,206 & 78,296 & 5.4h & 164 \\
Chicago & 17,291 & 38,671 & 118.0h & 92 \\
Austin & 5,968 & 6,029 & 23.8h & 117 \\
San Francisco & 17,723 & 17,007 & 15.1h & 37 \\
\bottomrule
\end{tabular}
\caption{the four feeds, two complete weeks from each.}
\end{table*}

Tickets still open when I pulled the data are kept and counted as slow. All of
them are months past any threshold I use. Chicago closes most of its feed within
a second, so I drop those rows. They are records, not work.

The task is simple. At the moment a ticket opens, predict whether it will take
longer than the median to close. Recovered skill is the share of the achievable
margin over chance that an arm reaches.

I need a ceiling to measure against. A prior reads two columns, the service
category and the owning department. So I fit the conditional rate straight from
those same two columns. That tells me what the columns support. Comparing against
it separates a weak rule from weak features. It is one estimator and not a proved
upper bound, so a rule can beat it.

I report two kinds of interval. A prior is a function of the category. It gives
the same answer to every ticket in a category. Resampling tickets treats 17,007
San Francisco rows as independent when I really have 37. So I resample both ways.
The category interval is the one a claim about unseen work has to survive [1].

I pulled each catalogue with the query limited to category names and row counts.
No duration, no rate, no outcome. I wrote the rules from that alone, committed
them with a prediction, and only then ran the evaluation. I amended one San
Francisco file 27 seconds after committing it, to match on word boundaries
instead of substrings. That was before the city's data existed on disk.

The biggest risk here is that the result describes one person's reasoning. So
Austin and San Francisco each got two sets of rules under the same blind
condition. I wrote one. A language model wrote the other from the same catalogue.
New York is the contrast, with rules written while I could see the real rates.
Chicago is where I carried those rules.

\section*{Results}

\subsection*{What a blind prior recovers}

\begin{table*}[t]
\centering
\footnotesize
\begin{tabular}{llrrrrr}
\toprule
city & rules & AUC & by ticket & by category & ceiling & recovered \\
\midrule
Austin & mine & 0.494 & [0.481, 0.508] & [0.331, 0.649] & 0.883 & $-1\%$ \\
Austin & the model & 0.611 & [0.598, 0.626] & [0.442, 0.770] & 0.883 & 29\% \\
San Francisco & mine & 0.676 & [0.670, 0.684] & [0.515, 0.815] & 0.886 & 46\% \\
San Francisco & the model & 0.777 & [0.770, 0.784] & [0.529, 0.859] & 0.886 & 72\% \\
\bottomrule
\end{tabular}
\caption{what blind rules recover, by city and by author. Ticket intervals are 1,000 resamples, category intervals 2,000.}
\end{table*}

The answer is not yes or no (Table 2). The best blind rules recover 72\% of the
achievable skill in San Francisco. The weakest recover nothing in Austin. Had I
tested one city I would have drawn a confident conclusion either way, and the
city would have decided which.

The two intervals disagree. By ticket, three of the four arms beat chance. By
category, only the two San Francisco arms do. Austin's model-written set spans
[0.442, 0.770], and I cannot tell it from guessing. The 72\% spans roughly 7\% to
93\% on the same unit. Read every claim below at that width.

\subsection*{What the rules add over structure alone}

A rule set can score well for two reasons. It may be reasoning. Or it may just be
cutting a lopsided taxonomy in any sensible way. Two controls tell them apart.

\begin{table*}[t]
\centering
\footnotesize
\begin{tabular}{llrrr}
\toprule
city & arm & AUC & [95\% CI] & recovered \\
\midrule
Austin & blind rules & 0.611 &  & 29\% \\
Austin & category frequency alone & 0.516 & [0.351, 0.726] & 4\% \\
Austin & same rules, verdicts shuffled & 0.501 & [0.322, 0.674] & 0\% \\
San Francisco & blind rules & 0.777 &  & 72\% \\
San Francisco & category frequency alone & 0.672 & [0.441, 0.785] & 45\% \\
San Francisco & same rules, verdicts shuffled & 0.499 & [0.312, 0.690] & 0\% \\
\bottomrule
\end{tabular}
\caption{blind rules against two baselines built from structure alone. Category intervals.}
\end{table*}

The shuffle keeps everything structural. Same taxonomy, same row counts, same set
of scores. It only breaks which category got which score. It lands at 0.501 and
0.499 (Table 3). So the reasoning is doing the work, not the shape of the
problem.

Frequency is the harder baseline. The pull gave me row counts, and row counts are
real operational data. Scored rarer-is-slower, frequency alone reaches 0.672 in
San Francisco. My own rules reach 0.676. My rules add almost nothing beyond
volume. The model's rules clear frequency by 0.105. Picking that direction is one
bit the catalogue did not give me, so the baseline is if anything generous.

\subsection*{The author effect is not established}

The model beats me by 0.116 in Austin and 0.101 in San Francisco. Same direction,
similar size, twice. That looks like something. By ticket both differences miss
zero, [0.104, 0.131] and [0.095, 0.107]. By category neither does, [-0.036,
0.288] and [-0.021, 0.174]. Two cities are not enough. An earlier draft of mine
reported these with no interval at all.

\subsection*{What I predicted, before I looked}

\begin{center}
\footnotesize
\setlength{\tabcolsep}{4pt}
\begin{tabular}{llrl}
\toprule
arm & predicted & observed & in band \\
\midrule
Austin, mine & 0.60 to 0.72 & 0.494 & no \\
Austin, the model & 0.52 to 0.62 & 0.611 & yes \\
SF, mine & 0.55 to 0.68 & 0.676 & yes \\
SF, the model & none & 0.777 & - \\
\bottomrule
\end{tabular}
\end{center}

I set the first band before the project had any result. It was wrong. I set the
next two after seeing Austin, so I already knew roughly what to expect. That is a
weaker achievement and I count it as one. The fourth arm produces my headline 72\%
and carries no prediction at all. The only prediction I made in real ignorance is
the one that failed.

\subsection*{One department explains most of the gap}

\begin{table*}[t]
\centering
\footnotesize
\begin{tabular}{llrrrr}
\toprule
subset & rules & n & prior & ceiling & recovered \\
\midrule
Austin, all tickets & mine & 6,029 & 0.494 & 0.883 & $-1\%$ \\
Austin, excluding ARR & mine & 4,008 & 0.644 & 0.881 & 38\% \\
Austin, all tickets & the model & 6,029 & 0.611 & 0.883 & 29\% \\
Austin, excluding ARR & the model & 4,008 & 0.766 & 0.881 & 70\% \\
San Francisco, all tickets & mine & 17,007 & 0.676 & 0.886 & 46\% \\
San Francisco, all tickets & the model & 17,007 & 0.777 & 0.886 & 72\% \\
\bottomrule
\end{tabular}
\caption{removing everything one Austin department owns, for both sets of rules.}
\end{table*}

Austin Resource Recovery (Table 4) is 34\% of the city's volume. 70\% of its work
runs slow. Both sets of rules call it fast, mean prediction 0.37. Take it out and
the model's rules reach 70\% against San Francisco's 72\%. Mine reach 38\% against
46\%. The gap closes for both of us.

The reason is plain. A missed collection sits open until the next scheduled
route, a week away. The work is routine. The ticket is not. Nothing in the phrase
ARR - Compost tells you that.

Two caveats. ARR is a department prefix, not a category called waste. So removing
it also removes fast work that department owns. ARR - Dead Animal Collection runs
0.0\% slow and ARR - General 8.9\%. Both were predicted fast and both were right.
San Francisco has the same problem at a quarter of the size. Graffiti Public is
8.2\% of volume, 85.1\% slow, and priced at 0.37. This subset does not remove it.

\subsection*{The finding}

A blind prior recovers about 70\% of the achievable skill on work whose length
follows from the job. It recovers nothing on work whose clock is an
administrative cycle. What decides its value is how much of a client's volume
sits in the second group, and you cannot see that in their taxonomy.

The 70\% is the model's, in both cities. I reach 38\% and 46\% on the same
comparison. The pattern holds for both of us. The level does not.

\subsection*{What the real rates are worth}

\begin{table*}[t]
\centering
\footnotesize
\begin{tabular}{llrrrr}
\toprule
rules & tested on & threshold & AUC & ceiling & recovered \\
\midrule
New York & New York & 5.4h & 0.840 & 0.918 & 81\% \\
NY rules carried & Chicago & 118.0h & 0.496 & 0.717 & $-2\%$ \\
NY rules carried & Chicago & 5.4h & 0.581 & 0.717 & 37\% \\
\bottomrule
\end{tabular}
\caption{the same procedure with the site's rates in hand, and those rules carried to another city.}
\end{table*}

Rules written with a city's rates in hand reach 81\% there. That is above the 72\%
a blind author got in San Francisco, but not by much.

Carried to Chicago at Chicago's own median they score 0.496. I called that
actively misleading in an earlier draft. It does not hold up (Table 5). The two
medians are 5.4h and 118.0h, 22 times apart. The carried row asks the rules a
question nobody wrote them for. Held at New York's own threshold the same rules
reach 0.581, or 37\%. What is left is mechanical. The rules fire on 17.1\% of
Chicago rows, because they are New York acronyms.

\subsection*{Distribution fidelity misses all of this}

\begin{table*}[t]
\centering
\footnotesize
\begin{tabular}{llrrr}
\toprule
task & generator & AUC & recovered & impossible \\
\midrule
city median & real records & 0.918 & 100\% & 0.0\% \\
city median & marginal only & 0.511 & 3\% & 70.5\% \\
city median & marginals + pairwise & 0.917 & 100\% & 0.0\% \\
city median & empirical joint & 0.918 & 100\% & 0.0\% \\
within category & real records & 0.524 & 100\% & 0.0\% \\
within category & marginal only & 0.502 & 10\% & 65.7\% \\
within category & marginals + pairwise & 0.506 & 25\% & 0.0\% \\
within category & empirical joint & 0.524 & 100\% & 0.0\% \\
\bottomrule
\end{tabular}
\caption{training on generated New York records and testing on real ones. Ten seeds.}
\end{table*}

I fit each generator on real New York records, sample from it, train a model on
the sample, and test on real records. Train on synthetic, test on real [2]
(Table 6). Matching marginals and pairwise dependence is the bar PrivBayes and
the NIST winner MST fit directly [3, 4]. It is roughly what SDV's quality report
scores afterwards [5, 6]. That report is a diagnostic run on a generator's
output, and similarity and usefulness may not correlate [6].

This ladder tests my first recorded hypothesis, fixed in \texttt{HYPOTHESIS.md} before I
ran anything. I expected generated records at that fidelity to hold combinations
the real process could never produce. I expected accuracy to fall as those rose.
It failed. The reason is structural rather than empirical, and the work moved to
the narrower question this paper answers. That is why my stated hypothesis and my
pre-registration are not the same sentence.

The pairwise generator draws department conditional on category. It cannot emit a
pair it never saw. So 0.0\% impossible is an identity, not a measurement. The
learner is additive over one-hot columns and has no interaction terms. A wrong
combination could not have hurt it either. The design settled that outcome before
I touched any data.

The city-median task cannot separate the two orders for another reason. Its label
is nearly a function of one column, and a pairwise generator gets that for free.
Labelling each ticket against its own category's median takes that column away.
Then they separate. Pairwise keeps 25\% where the joint keeps 100\%. The ceiling
there is 0.524, so the task holds little signal and the ratio sits on a small
base. It is the only measurement here that tells low-order fidelity from the full
joint. A distribution-level check scores both generators as near-perfect.

\section*{Limits}

Four cities, one domain, one task. Municipal work may be unusually tied to
institutions. A 311 feed has three usable columns, no documents that have to
agree with each other, and no prices. It is a thin stand-in for the enterprise
data the question is really about.

Two blind cities, not four. New York is a contamination control and Chicago its
transfer target. Neither replicates anything. What replicates is the author
effect, and its category intervals cross zero.

The second author is a language model, and this is the biggest hole in my design.
These feeds are among the most widely mirrored public datasets there are. Blind
means the session showed it no durations. It does not mean the weights hold none.
Models have memorised popular tabular datasets and score better on ones they have
seen [12]. That is the exact failure this arm is open to, and it produces every
headline number. I did not record the prompt or the model version, so I cannot
rerun it.

I wrote the San Francisco rules knowing what Austin had shown. They are blind to
San Francisco's rates. I am not blind to the lesson.

The ARR result is a subset I chose after seeing which family my rules got wrong,
in one city. It fits the numbers. Nothing has tested it. New York also carries
known artefacts. One department closes most of its tickets at exactly midnight,
so its durations are partly administrative.

That fidelity and utility can come apart is already known. The two rank
generators differently [7, 8]. High fidelity sits alongside poor classification
performance [9]. Models trained on synthetic data without care do not carry back
to real data [10]. Whether the standard metrics measure the right thing is itself
argued over [11]. I include the ladder because it shows a distribution-level
check cannot see the effect I measured.

\section*{What it would take to answer the general question}

Make the rules generate. My arms compare rule sets against a conditional fitted
on real data. That measures how much of a conditional a person can guess. The
question asks what generated data costs. To answer it I would sample rows from
the taxonomy and the volume counts, label them with the prior, train on those,
and test on real records. That turns a share of a margin into the number a buyer
needs. How much of the training value of real records does a generator built
without them deliver.

Predict the failures first. My ARR explanation says any category whose clock is
an administrative cycle will break a blind prior. That is checkable. Name those
categories in a fifth city from the taxonomy, write the list down, then measure.

More authors would help, including people who do the work. So would a language
model whose prompt and version I recorded.

One thing holds either way. You can only evaluate a prior as a prior once per
dataset, before anyone has seen the outcomes, and the prediction has to come
first. Of the predictions I made here, the only one made in real ignorance is the
one that turned out wrong.

\section*{References}

[1] C. A. Field and A. H. Welsh. Bootstrapping clustered data. \textit{J. R. Stat. Soc. B}, 69(3):369-390, 2007.

[2] C. Esteban, S. L. Hyland and G. Rätsch. Real-valued (medical) time series
generation with recurrent conditional GANs. arXiv:1706.02633, 2017.

[3] J. Zhang, G. Cormode, C. M. Procopiuc, D. Srivastava and X. Xiao. PrivBayes:
private data release via Bayesian networks. \textit{ACM Trans. Database Syst.},
42(4):article 25, 2017.

[4] R. McKenna, G. Miklau and D. Sheldon. Winning the NIST contest: a scalable
and general approach to differentially private synthetic data. \textit{J. Privacy and Confidentiality}, 11(3), 2021.

[5] N. Patki, R. Wedge and K. Veeramachaneni. The synthetic data vault. In \textit{IEEE DSAA}, pages 399-410, 2016.

[6] DataCebo. SDMetrics quality report: column shapes and column pair trends.
Software and documentation, doi:10.5281/zenodo.14279167, 2024.

[7] A. F. Karr, C. N. Kohnen, A. Oganian, J. P. Reiter and A. P. Sanil. A
framework for evaluating the utility of data altered to protect confidentiality.
\textit{The American Statistician}, 60(3):224-232, 2006.

[8] J. Snoke, G. M. Raab, B. Nowok, C. Dibben and A. Slavkovic. General and
specific utility measures for synthetic data. \textit{J. R. Stat. Soc. A},
181(3):663-688, 2018.

[9] L. Hansen, N. Seedat, M. van der Schaar and A. Petrovic. Reimagining
synthetic tabular data generation through data-centric AI: a comprehensive
benchmark. In \textit{NeurIPS Datasets and Benchmarks}, 2023.

[10] B. van Breugel, Z. Qian and M. van der Schaar. Synthetic data, real errors:
how (not) to publish and use synthetic data. In \textit{ICML}, PMLR 202:34793-34808,
2023.

[11] Y. Du and N. Li. Systematic assessment of tabular data synthesis. In \textit{ACM CCS}, 2025.

[12] S. Bordt, H. Nori, V. Rodrigues, B. Nushi and R. Caruana. Elephants never
forget: memorization and learning of tabular data in large language models. In
\textit{COLM}, 2024.

\end{document}
